\chapter*{Chapitre 7 — Régulation et GTB}
\addcontentsline{toc}{chapter}{Chapitre 7 — Régulation et GTB}

% ============================================================
\section*{Énoncés}
% ============================================================

\exercice{Choix de la loi de régulation}

Une chaudière à condensation dessert un réseau de radiateurs en acier. La consigne de température ambiante est $w = \SI{19}{\celsius}$. Après un arrêt de nuit, la mesure de température ambiante est $y = \SI{14}{\celsius}$.

\begin{enumerate}
  \item Calculer l'erreur initiale $e_0$.
  \item Avec un régulateur proportionnel de gain $K_P = \SI{5}{\percent\per\kelvin}$, calculer le signal de commande initial $u(t=0)$ sachant que $u_0 = 50\,\%$.
  \item La commande calculée est-elle physiquement réalisable ? Conclure sur la saturation de l'actionneur.
  \item Justifier l'intérêt d'ajouter une action intégrale (régulateur PI) pour atteindre exactement $\SI{19}{\celsius}$ en régime permanent.
\end{enumerate}

\exercice{Paramétrage de la loi d'eau d'une installation}

Un immeuble de logements à Lyon (zone H1b, $T_{ext,base} = \SI{-7}{\celsius}$) est chauffé par des radiateurs à basse température ($T_{dep,max} = \SI{55}{\celsius}$ à $T_{ext,base}$). Le chauffage se coupe à $T_{ext} = \SI{16}{\celsius}$ ($T_{dep,min} = \SI{30}{\celsius}$).

\begin{enumerate}
  \item Calculer la pente de la loi d'eau $K_{loi}$.
  \item Déterminer la température de départ programmée pour $T_{ext} = \SI{2}{\celsius}$.
  \item Un correcteur de confort ajoute $+\SI{3}{\kelvin}$ à la consigne de départ lorsque les occupants signalent une sensation de froid. Quelle est la nouvelle température de départ pour $T_{ext} = \SI{2}{\celsius}$ ?
  \item Quel est l'impact de cette correction sur la consommation de la chaudière ? (réponse qualitative)
\end{enumerate}

\exercice{Identification des composants GTB}

Un bâtiment tertiaire est équipé d'une GTB structurée en trois niveaux. Le technicien relève les équipements suivants :

\begin{enumerate}[label=\alph*)]
  \item Sonde de température ambiante PT1000 installée dans le bureau
  \item Logiciel SCADA affichant les synoptiques sur écran tactile au poste de garde
  \item Régulateur de zone Siemens Synco 700 gérant la vanne de la batterie chaude
  \item Vanne 2 voies motorisée DN15 sur circuit eau chaude
  \item Switch Ethernet administrable reliant les automates et le serveur
  \item Variateur de fréquence Schneider sur moteur pompe de circulation
\end{enumerate}

Classer chacun de ces équipements dans le niveau GTB correspondant (niveau 1, 2 ou 3) et préciser son rôle dans la chaîne de régulation.

% ============================================================
% EXERCICE 4
% ============================================================

\exercice{Régulation de pression différentielle par variateur de fréquence}

Un réseau hydraulique de chauffage alimente 8 radiateurs en bitube. La pompe de circulation est entraînée par un variateur de fréquence (VFD). Les données nominales à \SI{50}{\hertz} sont :

\begin{center}
\begin{tabular}{ll}
  \toprule
  Grandeur & Valeur \\
  \midrule
  Débit nominal & $Q_n = \SI{3.5}{\cubic\metre\per\hour}$ \\
  HMT nominale  & $H_n = \SI{8}{mCE}$ \\
  Puissance absorbée & $P_n = \SI{180}{\watt}$ \\
  Consigne pression différentielle & $\Delta P_{cons} = \SI{4}{mCE}$ (mode $\Delta P$ proportionnel) \\
  \bottomrule
\end{tabular}
\end{center}

\begin{enumerate}
  \item En mi-saison, la charge thermique est de 40\,\%. En mode $\Delta P$ proportionnel, le débit est proportionnel à la charge. Calculer le débit résiduel $Q_1$.
  \item Calculer la fréquence d'alimentation $f_1$ du moteur (loi de similitude : $Q \propto n \propto f$).
  \item Calculer la puissance absorbée $P_1$ à cette fréquence (loi de similitude : $P \propto n^3$).
  \item Calculer l'économie d'énergie sur une saison de chauffage de 4\,000\,h si la mi-saison représente 60\,\% du temps (2\,400\,h) et le plein régime 40\,\% (1\,600\,h).
  \item Calculer l'économie financière annuelle (électricité à \SI{0.18}{\euro\per\kilo\watt\hour}).
\end{enumerate}

% ============================================================
% EXERCICE 5
% ============================================================

\exercice{Calcul d'économie GTB : passage de classe C à classe B}

Un immeuble de bureaux de \SI{2000}{\metre\squared} (SHON) consomme actuellement \SI{85}{kWh$_\text{ef}$/m²/an} en chauffage, climatisation et ECS. Il est classé GTB\,C (régulation basique). Une rénovation GTB vers la classe\,B est envisagée.

\begin{center}
\begin{tabular}{ll}
  \toprule
  Donnée & Valeur \\
  \midrule
  Surface SHON & \SI{2000}{\metre\squared} \\
  Consommation actuelle (classe\,C) & 85 kWh$_\text{ef}$/m²/an \\
  Économie attendue C$\to$B (NF EN 15232) & 18\,\% \\
  Répartition & 65\,\% chauffage gaz, 35\,\% froid+auxiliaires élec \\
  Prix gaz & \SI{0.095}{\euro\per\kilo\watt\hour} \\
  Prix électricité & \SI{0.18}{\euro\per\kilo\watt\hour} \\
  Coût rénovation GTB & \SI{45000}{\euro} \\
  Durée de vie GTB & 15 ans \\
  \bottomrule
\end{tabular}
\end{center}

\begin{enumerate}
  \item Calculer la consommation totale annuelle actuelle en kWh$_\text{ef}$.
  \item Calculer la consommation économisée en kWh$_\text{ef}$ après passage en classe\,B.
  \item Calculer la répartition gaz/élec de l'économie et l'économie financière annuelle en euros.
  \item Calculer le TRI simple (Temps de Retour sur Investissement) de l'investissement GTB.
  \item Calculer le bénéfice total sur 15 ans en tenant compte d'une hausse des prix énergie de 3\,\%/an (somme géométrique).
  \item La RE2020 impose le décret BACS pour les bâtiments tertiaires au-delà de \SI{290}{\kilo\watt}. Quelle classe GTB est imposée au minimum ?
\end{enumerate}

\newpage
% ============================================================
\section*{Corrigés}
% ============================================================

\subsection*{Corrigé — Exercice 1}

\begin{enumerate}
  \item Erreur initiale :
  \[
    e_0 = w - y = 19 - 14 = \boxed{\SI{5}{\kelvin}}
  \]

  \item Signal de commande :
  \[
    u = K_P \cdot e_0 + u_0 = 5 \times 5 + 50 = \boxed{75\,\%}
  \]
  La vanne s'ouvre à $75\,\%$ de sa course maximale.

  \item La commande $75\,\%$ est physiquement réalisable (comprise entre 0 et 100\,\%). La saturation n'est pas atteinte dans ce cas. En revanche, pour $e > \SI{10}{\kelvin}$ : $u = 5 \times 10 + 50 = 100\,\%$ — saturation atteinte, la vanne est pleinement ouverte.

  \item En régime permanent, un actionneur P pur laisse subsister un écart statique résiduel. L'action intégrale accumule l'erreur dans le temps et ajuste progressivement $u$ jusqu'à ce que $e = 0$, garantissant ainsi $y = w = \SI{19}{\celsius}$ en régime établi, quelle que soit la charge thermique.
\end{enumerate}

\subsection*{Corrigé — Exercice 2}

\begin{enumerate}
  \item Pente de la loi d'eau :
  \[
    K_{loi} = \frac{T_{dep,max} - T_{dep,min}}{T_{ext,coupure} - T_{ext,base}} = \frac{55 - 30}{16 - (-7)} = \frac{25}{23} \approx \boxed{1{,}09}
  \]

  \item Température de départ à $T_{ext} = \SI{2}{\celsius}$ :
  \[
    T_{dep} = T_{dep,min} + K_{loi} \times (T_{ext,coupure} - T_{ext}) = 30 + 1{,}09 \times (16 - 2) = 30 + 15{,}3 \approx \boxed{\SI{45}{\celsius}}
  \]

  \item Avec correcteur de confort $+\SI{3}{\kelvin}$ :
  \[
    T_{dep,corr} = 45 + 3 = \boxed{\SI{48}{\celsius}}
  \]

  \item Une hausse de $\SI{3}{\kelvin}$ sur la température de départ augmente les pertes thermiques du réseau de distribution et réduit le rendement de condensation de la chaudière (la température de retour augmente). On estime une surconsommation de l'ordre de 6 à 10\,\% selon l'isolation du réseau. Le correcteur doit donc être utilisé avec modération et remis à zéro dès retour à l'état de confort.
\end{enumerate}

\subsection*{Corrigé — Exercice 3}

\begin{center}
\begin{tabular}{clll}
  \toprule
  Équipement & Niveau & Rôle & Nature \\
  \midrule
  a) Sonde PT1000 & 1 — Terrain & Mesure grandeur réglée ($T_{amb}$) & Capteur \\
  b) Logiciel SCADA & 3 — Supervision & Visualisation, historique, alarmes & IHM \\
  c) Régulateur Synco 700 & 2 — Automatisme & Calcule la commande (boucle PI) & Contrôleur \\
  d) Vanne motorisée & 1 — Terrain & Agit sur le débit d'eau chaude & Actionneur \\
  e) Switch Ethernet & 3 — Supervision & Infrastructure réseau IP & Communication \\
  f) Variateur de fréquence & 1/2 — Terrain/Auto. & Adapte la vitesse de la pompe & Actionneur piloté \\
  \bottomrule
\end{tabular}
\end{center}

\begin{methode}
Note : le variateur de fréquence est physiquement au niveau terrain mais intègre souvent une intelligence locale (boucle PID de pression) — il peut donc être considéré comme niveau 1 ou niveau 2 selon le contexte.
\end{methode}

\subsection*{Corrigé — Exercice 4}

\begin{enumerate}
  \item \textbf{Débit résiduel à 40\,\% de charge :}

  En mode $\Delta P$ proportionnel, le débit est proportionnel à la charge thermique :
  \[
    Q_1 = 0{,}40 \times Q_n = 0{,}40 \times 3{,}5 = \boxed{\SI{1.4}{\cubic\metre\per\hour}}
  \]

  \item \textbf{Fréquence $f_1$ du moteur :}

  Par la loi de similitude ($Q \propto n \propto f$) :
  \[
    \frac{Q_1}{Q_n} = \frac{f_1}{f_n}
    \quad\Rightarrow\quad
    f_1 = f_n \times \frac{Q_1}{Q_n} = 50 \times \frac{1{,}4}{3{,}5} = \boxed{\SI{20}{\hertz}}
  \]

  \item \textbf{Puissance absorbée à mi-charge :}

  Par la loi de similitude ($P \propto n^3$) :
  \[
    P_1 = P_n \times \left(\frac{f_1}{f_n}\right)^3 = 180 \times \left(\frac{20}{50}\right)^3
         = 180 \times (0{,}40)^3 = 180 \times 0{,}064 = \boxed{\SI{11.5}{\watt}}
  \]

  \item \textbf{Économie d'énergie sur la saison :}

  Énergie consommée \textit{sans} VFD (pompe à vitesse fixe, $P_n$ permanente) :
  \[
    E_{sans} = P_n \times \text{durée totale} = \frac{180}{1000} \times 4\,000 = \SI{720}{\kilo\watt\hour}
  \]

  Énergie consommée \textit{avec} VFD :
  \begin{align*}
    E_{avec} &= P_n \times t_{plein} + P_1 \times t_{mi} \\
             &= \frac{180}{1000} \times 1\,600 + \frac{11{,}5}{1000} \times 2\,400 \\
             &= 288 + 27{,}6 = \SI{315.6}{\kilo\watt\hour}
  \end{align*}

  Économie d'énergie :
  \[
    \Delta E = E_{sans} - E_{avec} = 720 - 315{,}6 = \boxed{\SI{404.4}{\kilo\watt\hour/saison}}
  \]

  Soit une réduction de $\dfrac{404{,}4}{720} \approx \mathbf{56\,\%}$ de la consommation de la pompe.

  \item \textbf{Économie financière annuelle :}
  \[
    \text{Économie} = \Delta E \times p_{elec} = 404{,}4 \times 0{,}18 = \boxed{\SI{72.8}{\euro/an}}
  \]
\end{enumerate}

\begin{methode}
Les lois de similitude des turbomachines (lois d'affinité) sont fondamentales en CVC. Retenir : $Q \propto n$, $H \propto n^2$, $P \propto n^3$. La réduction de $P$ est cubique : à demi-vitesse, la puissance est divisée par 8.
\end{methode}

\subsection*{Corrigé — Exercice 5}

\begin{enumerate}
  \item \textbf{Consommation totale annuelle actuelle :}
  \[
    C_{tot} = S \times c_{sp} = 2\,000 \times 85 = \boxed{170\,000\;\text{kWh}_\text{ef}/\text{an}}
  \]

  \item \textbf{Économie après passage en classe\,B :}
  \[
    \Delta C = C_{tot} \times 18\,\% = 170\,000 \times 0{,}18 = \boxed{30\,600\;\text{kWh}_\text{ef}/\text{an}}
  \]

  \item \textbf{Répartition gaz / électricité et économie financière :}

  Économie sur la part gaz (65\,\%) :
  \[
    \Delta C_{gaz} = 30\,600 \times 0{,}65 = 19\,890\;\text{kWh/an}
    \quad\Rightarrow\quad
    \text{Éco}_{\text{gaz}} = 19\,890 \times 0{,}095 = \SI{1889.6}{\euro/an}
  \]

  Économie sur la part électricité (35\,\%) :
  \[
    \Delta C_{elec} = 30\,600 \times 0{,}35 = 10\,710\;\text{kWh/an}
    \quad\Rightarrow\quad
    \text{Éco}_{\text{elec}} = 10\,710 \times 0{,}18 = \SI{1927.8}{\euro/an}
  \]

  \textbf{Économie financière annuelle totale :}
  \[
    \text{Éco}_{tot} = 1\,889{,}6 + 1\,927{,}8 = \boxed{\SI{3817.4}{\euro/an}}
  \]

  \item \textbf{Temps de Retour sur Investissement simple (TRI) :}
  \[
    TRI = \frac{C_{investissement}}{\text{Éco}_{tot}} = \frac{45\,000}{3\,817{,}4} \approx \boxed{11{,}8\;\text{ans}}
  \]

  \item \textbf{Bénéfice total sur 15 ans avec hausse énergétique 3\,\%/an :}

  La somme géométrique des économies revalorisées (hausse de 3\,\%/an) est :
  \[
    B_{15} = \text{Éco}_{tot} \times \sum_{k=0}^{14}(1{,}03)^k
           = \text{Éco}_{tot} \times \frac{(1{,}03)^{15} - 1}{1{,}03 - 1}
  \]

  Calcul intermédiaire :
  \[
    (1{,}03)^{15} \approx 1{,}5580
    \quad\Rightarrow\quad
    \frac{1{,}5580 - 1}{0{,}03} = \frac{0{,}5580}{0{,}03} \approx 18{,}60
  \]

  \[
    B_{15} = 3\,817{,}4 \times 18{,}60 \approx \boxed{\SI{71003}{\euro}}
  \]

  Bénéfice net sur 15 ans (après déduction de l'investissement initial) :
  \[
    B_{net} = 71\,003 - 45\,000 = \boxed{\SI{26003}{\euro}}
  \]

  \item \textbf{Classe GTB imposée par le décret BACS :}

  Le décret BACS (n°2020-658) impose l'installation d'un système d'automatisation et de contrôle pour les bâtiments tertiaires dont la puissance nominale des systèmes CVC dépasse \SI{290}{\kilo\watt}. La classe minimale requise est la \textbf{classe\,B} au sens de la norme NF EN 15232 / EN\,52120-1, correspondant à une automatisation avancée (régulation PI, loi d'eau, optimisation du démarrage, suivi énergétique).

  \begin{attention}
    Le décret BACS ne mentionne pas explicitement les classes A–D de la norme EN\,15232, mais les fonctionnalités qu'il impose (monitoring, journalisation, pilotage centralisé) correspondent au minimum à la \textbf{classe\,B}. Le passage vers la classe\,A (GTB intégrée multi-énergie) est fortement recommandé pour les bâtiments soumis au Décret Tertiaire (Éco Énergie Tertiaire).
  \end{attention}
\end{enumerate}
